\section{Discusión preliminar y limitaciones}
El aporte potencial no consiste únicamente en incorporar interfaces, sino en especificar un modelo reproducible de seguimiento programático, control autorizado de cambios y trazabilidad. La inspección del código proporciona evidencia técnica del antecedente, pero no prueba el beneficio de PID 2.0.

No se han ejecutado entrevistas, pruebas de aceptación ni observación funcional. Permanecen pendientes la autoridad competente para aprobar cambios, la fórmula de avance, la identidad de actividades entre versiones y la evaluación de seguridad. En la sesión 3 se prevé resolver el modelo de datos y flujos; en la sesión 4, implementar y evaluar una PoC delimitada con datos sintéticos. No se formulan conclusiones finales.
